% =============================================================
% Chapitre II — État de l'art du downscaling par apprentissage profond
% Style cible : gabarit Donald FEUZING NTEMMA (URIFIA 2024)
% =============================================================

\documentclass[14pt,a4paper,openany]{extreport}

% ---------- Encodage et langue ----------
\usepackage[T1]{fontenc}
\usepackage[utf8]{inputenc}
\usepackage[french]{babel}

% ---------- Police lisible pour impression papier ----------
\usepackage{lmodern}

% ---------- Interligne (impression papier reliée) ----------
\usepackage{setspace}
\onehalfspacing

% ---------- Géométrie (marges adaptées à la reliure) ----------
\usepackage[a4paper,inner=3cm,outer=2.3cm,top=2.5cm,bottom=2.7cm,headheight=15pt,bindingoffset=0.3cm]{geometry}

% ---------- Mathématiques ----------
\usepackage{amsmath,amssymb,amsthm}

% ---------- Graphiques et couleurs ----------
\usepackage{graphicx}
\usepackage{xcolor}
\definecolor{uridarkblue}{RGB}{30,40,90}
\definecolor{urigray}{RGB}{70,70,70}

% ---------- Tableaux ----------
\usepackage{array}
\usepackage{booktabs}
\usepackage{multirow}

% ---------- Boîte SOMMAIRE et placeholders de figures ----------
\usepackage[most]{tcolorbox}

% ---------- Mise en forme des chapitres et sections ----------
\renewcommand{\thechapter}{\Roman{chapter}}
\renewcommand{\thesection}{\thechapter.\arabic{section}}
\renewcommand{\thesubsection}{\thesection.\arabic{subsection}}
\renewcommand{\thesubsubsection}{\thesubsection.\arabic{subsubsection}}

\makeatletter
\renewcommand{\@makechapterhead}[1]{%
  \vspace*{20\p@}%
  {\parindent \z@ \centering \normalfont
    \textsc{\large Chapitre}\par
    \vspace{0.3em}
    {\color{uridarkblue}\fontsize{60}{72}\selectfont\bfseries\thechapter\par}
    \vspace{1em}
    {\color{uridarkblue}\Large\bfseries\MakeUppercase{#1}\par}
    \vspace{1em}
    \rule{0.4\linewidth}{0.6pt}\par
    \vspace{2em}
  }}
\renewcommand{\@makeschapterhead}[1]{%
  \vspace*{20\p@}%
  {\parindent \z@ \centering \normalfont
    {\color{uridarkblue}\Large\bfseries\MakeUppercase{#1}\par}
    \vspace{1em}
    \rule{0.4\linewidth}{0.6pt}\par
    \vspace{2em}
  }}
\makeatother

\usepackage{titlesec}
\titleformat{\section}
  {\normalfont\Large\bfseries\color{uridarkblue}}
  {\thesection.}{0.8em}{}
\titleformat{\subsection}
  {\normalfont\large\bfseries}
  {\thesubsection.}{0.7em}{}
\titleformat{\subsubsection}
  {\normalfont\normalsize\bfseries\itshape}
  {\thesubsubsection.}{0.6em}{}

% ---------- En-têtes et pieds de page ----------
\usepackage{fancyhdr}
\pagestyle{fancy}
\fancyhf{}
\fancyhead[L]{\textsc{\nouppercase{\leftmark}}}
\fancyhead[R]{\thepage}
\fancyfoot[L]{\footnotesize Mémoire de Master en Informatique, Université de Dschang}
\fancyfoot[R]{\footnotesize URIFIA}
\renewcommand{\headrulewidth}{0.4pt}
\renewcommand{\footrulewidth}{0.2pt}
\renewcommand{\chaptermark}[1]{\markboth{#1}{}}
\renewcommand{\sectionmark}[1]{\markright{\thesection\ #1}{}}

% ---------- Légendes ----------
\usepackage[font=small,labelfont={sc,bf},textfont=it]{caption}
\renewcommand{\thefigure}{\arabic{figure}}
\renewcommand{\thetable}{\arabic{table}}

% Pour numérotation continue avec Chap I — démarrage à 9 figures, 1 tableau
\setcounter{figure}{8}
\setcounter{table}{0}

% ---------- Citations numériques ----------
\usepackage[numbers,square,sort&compress]{natbib}

% ---------- Hyperref ----------
\usepackage[hidelinks]{hyperref}

% ---------- Macros ----------
\newcommand{\Oracle}{\textsc{Oracle}}
\newcommand{\CorrDiff}{\textsc{CorrDiff}}
\newcommand{\LR}{\ensuremath{\text{LR}}}
\newcommand{\HR}{\ensuremath{\text{HR}}}
\newcommand{\xHR}{\ensuremath{x_{\HR}}}
\newcommand{\yLR}{\ensuremath{y_{\LR}}}
\newcommand{\muHR}{\ensuremath{\mu_{\HR}}}

\newtcolorbox{sommairebox}{
  enhanced,colback=white,colframe=urigray,boxrule=0.5pt,arc=0pt,
  left=1.5em,right=1.5em,top=0.8em,bottom=0.8em,
  title={\textsc{Sommaire}},
  attach boxed title to top center={yshift=-2mm,xshift=0mm},
  boxed title style={colback=white,colframe=white,coltitle=urigray,boxrule=0pt,
    left=0.8em,right=0.8em,top=0pt,bottom=0pt},
  fonttitle=\bfseries\small,
}

\newcommand{\figplaceholder}[2][0.9\linewidth]{%
  \begin{tcolorbox}[width=#1,height=5cm,valign=center,colback=gray!8,colframe=gray!45,arc=2pt,boxrule=0.5pt]
    \centering\small\textit{#2}
  \end{tcolorbox}%
}

\setcounter{chapter}{1}
\setcounter{page}{29}

\begin{document}

\chapter{État de l'art du downscaling par apprentissage profond}
\label{chap:etatdelart}

\begin{sommairebox}
\small
\noindent II.1 \ Introduction \dotfill \pageref{sec:soa-intro}\\
\noindent II.2 \ Approches déterministes \dotfill \pageref{sec:soa-det}\\
\noindent II.3 \ Approches probabilistes adversariales \dotfill \pageref{sec:soa-adv}\\
\noindent II.4 \ Approches à vraisemblance explicite \dotfill \pageref{sec:soa-vrai}\\
\noindent II.5 \ Modèles de diffusion \dotfill \pageref{sec:soa-diff}\\
\noindent II.6 \ Le paradigme à deux étages : CorrDiff \dotfill \pageref{sec:soa-corrdiff}\\
\noindent II.7 \ Fonctions de perte et contraintes physiques \dotfill \pageref{sec:soa-pertes}\\
\noindent II.8 \ Cadres d'évaluation et foundation models \dotfill \pageref{sec:soa-benchmarks}\\
\noindent II.9 \ Synthèse, tableau comparatif et gap scientifique \dotfill \pageref{sec:soa-synthese}
\end{sommairebox}

\vspace{1.5em}

% =============================================================
\section{Introduction}
\label{sec:soa-intro}
% =============================================================

Le chapitre précédent a identifié un trilemme pour le downscaling climatique génératif (section~I.7) : stabilité d'entraînement, réalisme des extrêmes et robustesse sous changement de distribution apparaissent inconciliables dans les approches existantes. Le présent chapitre confronte cette analyse au paysage effectif de la littérature, des CNN déterministes aux modèles de diffusion à deux étages. Nous abordons successivement les approches déterministes (section~\ref{sec:soa-det}), adversariales (section~\ref{sec:soa-adv}), à vraisemblance explicite (section~\ref{sec:soa-vrai}) et de diffusion (section~\ref{sec:soa-diff}), avant de détailler le paradigme \CorrDiff{} (section~\ref{sec:soa-corrdiff}) dont notre proposition est la causalisation. Pertes (section~\ref{sec:soa-pertes}) et cadres d'évaluation (section~\ref{sec:soa-benchmarks}) complètent la revue ; la synthèse (section~\ref{sec:soa-synthese}) identifie le gap qui motive notre contribution.

\begin{figure}[ht]
  \centering
  \resizebox{\linewidth}{!}{%
  \begin{tikzpicture}[
    yscale=1.0, xscale=0.92,
    timearrow/.style={-{Stealth[length=2.5mm]}, very thick, uridarkblue},
    tick/.style={uridarkblue, thick},
    year/.style={font=\scriptsize\bfseries, text=uridarkblue, anchor=north},
    detlabel/.style={font=\scriptsize, anchor=south, text=urigray, align=center,
                     rectangle, draw=urigray!50, fill=urigray!8, rounded corners=1pt,
                     inner sep=2pt, minimum height=5mm},
    genlabel/.style={font=\scriptsize, anchor=south, text=uriaccent!85!black, align=center,
                     rectangle, draw=uriaccent!60, fill=uriaccent!10, rounded corners=1pt,
                     inner sep=2pt, minimum height=5mm},
    diflabel/.style={font=\scriptsize, anchor=north, text=uridarkblue, align=center,
                     rectangle, draw=uridarkblue!60, fill=uridarkblue!10, rounded corners=1pt,
                     inner sep=2pt, minimum height=5mm},
    paradigm/.style={font=\scriptsize\itshape, text=urigray},
  ]
    % --- Timeline arrow ---
    \draw[timearrow] (-0.4, 0) -- (14.2, 0);

    % --- Year ticks ---
    \foreach \x/\y in {0/2014, 1.8/2017, 3.6/2020, 5.4/2022, 7.2/2023, 9.0/2024, 10.8/2025, 13.0/2026} {
      \draw[tick] (\x, 0.18) -- (\x, -0.18);
      \node[year] at (\x, -0.28) {\y};
    }

    % --- Determinist contributions (au-dessus, gris) ---
    \node[detlabel] (srcnn)  at (0.0, 0.4)  {SRCNN \\ \scriptsize Dong+'14};
    \node[detlabel] (deepsd) at (1.8, 0.4)  {DeepSD \\ \scriptsize Vandal+'17};

    % --- Generative contributions (au-dessus, accent rouge) ---
    \node[genlabel] (ddpm)    at (3.6, 1.5) {DDPM \\ \scriptsize Ho+'20};
    \node[genlabel] (edm)     at (5.4, 1.5) {EDM \\ \scriptsize Karras+'22};
    \node[genlabel] (sr3)     at (5.4, 0.4) {SR3 \\ \scriptsize Saharia+'22};
    \node[genlabel] (corrdiff)at (7.2, 1.5) {CorrDiff \\ \scriptsize Mardani+'23};
    \node[genlabel] (climax)  at (7.2, 0.4) {ClimaX / \\ GraphCast '23};
    \node[genlabel] (gencast) at (9.0, 1.5) {GenCast \\ \scriptsize Price+'24};
    \node[genlabel] (rampal)  at (9.0, 0.4) {cGAN-NZ \\ \scriptsize Rampal+'24};

    % --- Oracle (en-dessous, accent fort) ---
    \node[diflabel, fill=uriaccent!25, draw=uriaccent, very thick] (oracle)
      at (13.0, -0.55) {\Oracle{} \\ \scriptsize ce mémoire};

    % --- Paradigms (sous l'axe) ---
    \draw[urigray, dashed, thick] (-0.4,-1.4) -- (3.0,-1.4)
      node[midway, paradigm, below, yshift=-1pt] {\textbf{Déterministe}};
    \draw[uriaccent!70, thick] (3.0,-1.4) -- (10.8,-1.4)
      node[midway, paradigm, below, text=uriaccent, yshift=-1pt] {\textbf{Génératif (single-stage puis two-stage)}};
    \draw[uridarkblue, very thick] (10.8,-1.4) -- (14.0,-1.4)
      node[midway, paradigm, below, text=uridarkblue, yshift=-1pt] {\textbf{Two-stage causal}};

    % Connector to Oracle
    \draw[uriaccent, very thick, dashed] (13.0, -0.18) -- (oracle.north);
  \end{tikzpicture}%
  }
  \caption{Frise chronologique des principales contributions du downscaling climatique profond. Le glissement progressif du paradigme déterministe vers le paradigme génératif à deux étages est visible à partir de 2020.}
  \label{fig:frise}
\end{figure}

La figure~\ref{fig:frise} situe chronologiquement les contributions examinées : glissement déterministe $\to$ génératif à partir de 2020, accélération après 2022 sous l'effet conjoint des modèles de diffusion et de la maturation des GCM haute résolution.

% =============================================================
\section{Approches déterministes}
\label{sec:soa-det}
% =============================================================

\subsection{Régression statistique classique}
\label{ssec:soa-reg-classique}

Les approches statistiques de première génération --- régression linéaire, \emph{Bias Correction} (BC), \emph{Quantile Mapping} (QM), perfect prog --- ont été présentées au chapitre précédent (section~\ref{sec:classiques}). Elles partagent trois limites structurelles : stationnarité de la relation prédicteurs-prédictands, absence de cohérence spatiale (traitement pixel par pixel), et absence de calibration probabiliste. D'où l'introduction de méthodes d'apprentissage non-linéaires structurées spatialement.

\subsection{Réseaux convolutifs et super-résolution}
\label{ssec:soa-cnn}

L'application des CNN au downscaling s'est inspirée des architectures de super-résolution d'image, d'abord SRCNN~\cite{dong_srcnn_2014} puis EDSR~\cite{lim_edsr_2017}. \textbf{DeepSD}~\cite{vandal_deepsd_2017} fut l'une des premières architectures CNN dédiées au downscaling de la précipitation, par empilement de convolutions à champ réceptif croissant. Le \textbf{U-Net}~\cite{ronneberger_unet_2015} a connu une adoption massive, particulièrement sous sa forme conditionnée par des champs de surface (élévation). D'autres variantes (\textbf{PAN}~\cite{pan_pan_2019}, modules d'attention spatiale, pyramides de caractéristiques) ont été explorées.

Ces approches partagent une limite décisive : produisant une unique sortie HR par entrée, elles approximent la moyenne conditionnelle de la distribution cible. Pour la précipitation à queue lourde, cette moyenne est biaisée à la baisse aux valeurs élevées, d'où le lissage des extrêmes documenté dans la littérature.

\subsection{Architectures déterministes avancées}
\label{ssec:soa-avancees}

La période 2021-2023 a vu émerger des architectures plus sophistiquées, regroupées ici car elles partagent la même limite probabiliste que les CNN. Les \textbf{opérateurs neuronaux de Fourier} (FNO)~\cite{li_fno_2021} et leur variante sphérique \textbf{SFNO}~\cite{bonev_sfno_2023} apprennent des opérateurs entre fonctions par représentation spectrale, prometteurs en prévision globale mais expérimentaux en downscaling régional. Les \textbf{Transformers de vision} appliqués à la super-résolution, notamment \textbf{SwinIR}~\cite{liang_swinir_2021} et son adaptation climat \textbf{PrecipFormer}~\cite{lopezgomez_precipformer_2023}, exploitent l'attention pour capter des dépendances longues. Toutes conservent la limite fondamentale du déterminisme : pas de distribution conditionnelle calibrée, donc le verrou « réalisme des extrêmes » reste non adressé.

% =============================================================
\section{Approches probabilistes adversariales}
\label{sec:soa-adv}
% =============================================================

\subsection{Réseaux antagonistes génératifs conditionnels}
\label{ssec:soa-cgan}

Le glissement vers le probabiliste a d'abord été opéré par les \textbf{réseaux antagonistes génératifs} (GAN)~\cite{goodfellow_gan_2014} et leur variante conditionnelle (cGAN)~\cite{mirza_cgan_2014}. Un générateur $G$ transforme un bruit $z \sim \mathcal{N}(0,I)$ et un conditionnement $\yLR$ en échantillon plausible $\hat{x}_{\HR} = G(\yLR, z)$ ; un discriminateur $D$ apprend à distinguer vrais et faux. L'entraînement min-max alterne ces objectifs opposés. Appliqué au downscaling, le cGAN promet un ensemble d'échantillons HR plausibles permettant en principe d'estimer la variabilité prédictive et la distribution des extrêmes.

\subsection{cGAN résiduel pour le downscaling climatique}
\label{ssec:soa-rampal}

Une architecture représentative est le \textbf{Residual cGAN} de Rampal et~al.~\cite{rampal_cgan_2024}, qui combine un U-Net déterministe pré-entraîné et un cGAN sur le résidu : le générateur reçoit la prédiction U-Net, l'élévation HR, le champ GCM et un bruit, et produit une carte résiduelle additionnée à la prédiction déterministe. L'application à la Nouvelle-Zélande (CMIP6 ACCESS-CM2) combine trois termes de coût : MSE, perte adversariale pondérée par $\lambda_{\mathrm{adv}}$, et contrainte d'intensité pénalisant la surestimation des extrêmes.

D'autres travaux apparentés incluent Harris et~al.~\cite{harris_2022} (désagrégation de précipitations sur domaine européen) et Price et Rasp~\cite{price_rasp_2022} (prévision d'extrêmes à court terme via GAN).

\subsection{Limites empiriques des cGAN appliqués au downscaling}
\label{ssec:soa-cgan-limites}

Quatre limites empiriques caractérisent la famille adversariale. (i) \textbf{Forte sensibilité à $\lambda_{\mathrm{adv}}$} : la performance sur les extrêmes suit une courbe en cloche (trop faible $\Rightarrow$ lissage par la MSE ; trop élevé $\Rightarrow$ surestimation), avec un optimum étroit et peu transférable. (ii) \textbf{Sub-dispersion} : les analyses de Rampal et~al.~\cite{rampal_cgan_2024} et Mardani et~al.~\cite{mardani_corrdiff_2023} rapportent $\sigma_{\mathrm{pred}}/\mathrm{RMSE} < 0{,}5$, ensembles systématiquement trop confiants. (iii) \textbf{Risque de \emph{mode collapse}}, intrinsèque au cadre adversarial, latent même s'il n'est pas systématiquement observé en pratique. (iv) \textbf{Validité restreinte du cadre d'évaluation} : tests menés majoritairement sur couples synthétiques RCM-LR/RCM-HR (cadre « parfait »), la transférabilité aux conditions opérationnelles GCM brut restant incertaine.

% =============================================================
\section{Approches à vraisemblance explicite}
\label{sec:soa-vrai}
% =============================================================

Les \textbf{auto-encodeurs variationnels} (VAE)~\cite{kingma_vae_2014} apprennent une représentation latente probabiliste en maximisant la borne inférieure de la log-vraisemblance (ELBO). Appliqués au downscaling~\cite{lange_vae_2022}, ils produisent des sorties probabilistes mais souffrent d'un excès de lissage lié à l'ELBO et à la dimension réduite de l'espace latent. Les \textbf{flots de normalisation} appliquent une séquence de transformations inversibles différentiables à un bruit gaussien et autorisent l'évaluation exacte de la vraisemblance via le déterminant jacobien. Leur application au downscaling~\cite{groenke_flow_2020} a fourni des résultats compétitifs, mais l'expressivité reste bornée par la contrainte d'inversibilité.

Ces deux familles n'atteignent pas le niveau des GAN ou de la diffusion en downscaling climatique. Leur intérêt méthodologique --- évaluation exacte de la vraisemblance --- ne compense pas une expressivité insuffisante face aux extrêmes à queue lourde et à l'hétérogénéité spatiale de la précipitation fine. Elles restent à ce jour davantage des objets d'étude méthodologique que des candidats sérieux au downscaling opérationnel.\label{ssec:soa-vae}\label{ssec:soa-flow}\label{ssec:soa-vrai-bilan}

% =============================================================
\section{Modèles de diffusion}
\label{sec:soa-diff}
% =============================================================

\subsection{Denoising Diffusion Probabilistic Models}
\label{ssec:soa-ddpm}

Les \textbf{Denoising Diffusion Probabilistic Models} (DDPM)~\cite{ho_ddpm_2020} définissent un processus direct qui bruite progressivement le signal et apprennent à inverser ce processus pour reconstruire des échantillons. Le processus direct est une chaîne de Markov gaussienne :
\begin{equation}
  q(x_t \mid x_{t-1}) \;=\; \mathcal{N}\!\left(x_t;\, \sqrt{1 - \beta_t}\, x_{t-1},\, \beta_t I\right),
  \label{eq:ddpm-forward}
\end{equation}
avec un calendrier de variance $(\beta_t)_{t=1}^{T}$. L'objectif simplifié réduit l'apprentissage à la prédiction du bruit :
\begin{equation}
  \mathcal{L}_{\mathrm{simple}}(\theta) \;=\; \mathbb{E}_{t, x_0, \epsilon}\!\left[\,\left\| \epsilon - \epsilon_\theta(x_t, t) \right\|^2\,\right],
  \label{eq:ddpm-loss}
\end{equation}
où $\epsilon_\theta$ est typiquement un U-Net. L'échantillonnage inverse part d'un bruit pur. Face aux GAN, DDPM offre trois avantages pour le downscaling : entraînement stable sans déséquilibre adversarial, expressivité multimodale sans collapse, et calibration probabiliste intrinsèque liée à la vraisemblance variationnelle.

\subsection{Formalisme EDM de Karras et al.}
\label{ssec:soa-edm}

Le formalisme EDM~\cite{karras_edm_2022} reformule la diffusion dans un cadre théorique général. Sa contribution principale est un \textbf{préconditionnement systématique} de l'entrée et de la sortie via quatre fonctions $c_{\mathrm{in}}, c_{\mathrm{out}}, c_{\mathrm{noise}}, c_{\mathrm{skip}}$ dépendant de $\sigma$, qui unifie DDPM/score matching/VP/VE et stabilise l'entraînement. Le sampler de Heun à pas variable atteint une qualité comparable à DDPM avec moitié moins d'étapes. La calibration de $\sigma_{\mathrm{data}}$ (écart-type des données normalisées) est cruciale.

\subsection{Premières applications au climat : SR3, ResShift, ResDiff}
\label{ssec:soa-sr3-resdiff}

\textbf{SR3}~\cite{saharia_sr3_2022} (\emph{Super-Resolution via Repeated Refinement}) fut la première application notable de la diffusion conditionnelle à la super-résolution d'image, essaimant ensuite vers le climat. \textbf{ResShift}~\cite{yue_resshift_2023} accélère la diffusion par déplacement résiduel (diffusion partant de l'image LR sur-échantillonnée, corrigeant un résidu en peu d'étapes), avec une proximité manifeste vers le paradigme à deux étages de la section~\ref{sec:soa-corrdiff}. Pour mémoire, la version preprint de Mardani et~al.~\cite{mardani_corrdiff_2023} fut initialement nommée \emph{Residual Corrective Diffusion} (ResDiff) avant d'être renommée CorrDiff. Plusieurs évolutions récentes (depuis 2023), non encore appliquées au downscaling, pourraient à terme remplacer l'étage de diffusion : \textbf{flow matching}~\cite{lipman_flow_2023} (champ de vecteurs de transport, échantillonnage ODE), \textbf{consistency models}~\cite{song_consistency_2023} (échantillonnage en une étape), \textbf{diffusion latente}~\cite{rombach_latent_2022} (espace compressé). La diffusion EDM demeure notre choix de référence.\label{ssec:soa-flow-matching}

% =============================================================
\section{Le paradigme à deux étages : CorrDiff}
\label{sec:soa-corrdiff}
% =============================================================

\subsection{Idée fondatrice}
\label{ssec:soa-corrdiff-idee}

Le modèle \CorrDiff{} (\emph{Corrective Diffusion}) de Mardani et~al.~\cite{mardani_corrdiff_2023} (NVIDIA Research) propose un paradigme à deux étages dont la décomposition formelle est :
\begin{equation}
  \xHR \;=\; \mathrm{baseline}_{\LR \uparrow} \;+\; \muHR\!\left(\yLR, s_{\HR}\right) \;+\; \delta\!\left(\,\cdot\,\right),
  \label{eq:corrdiff}
\end{equation}
où $\mathrm{baseline}_{\LR \uparrow}$ désigne un sur-échantillonnage bilinéaire, $\muHR$ la moyenne HR apprise par régression déterministe, et $\delta$ un résiduel stochastique échantillonné par diffusion conditionnée sur $\muHR$ et $\yLR$. L'application visée est le downscaling kilométrique de sorties GCM avec un facteur de réduction de 25 à 50, le terrain initial étant Taïwan.

\subsection{Architecture détaillée}
\label{ssec:soa-corrdiff-archi}

Le premier étage est un U-Net classique entraîné par MSE entre $\muHR$ et la cible HR. Le second est un U-Net de diffusion EDM entraîné à prédire $\delta = \xHR - (\mathrm{baseline}_{\LR \uparrow} + \muHR)$ conditionnellement à $\muHR$ et $\yLR$, avec échantillonnage de Heun à 20-50 pas.

\begin{figure}[ht]
  \centering
  \figplaceholder{Figure F4 (Chap II) à produire --- architecture détaillée CorrDiff : (Étage 1) Entrée LR + élévation $\to$ U-Net régression $\to$ moyenne $\mu_{HR}$ ; (Étage 2) bruit + $\mu_{HR}$ + LR $\to$ U-Net diffusion EDM (préconditionnement $c_{in}, c_{out}, c_{noise}, c_{skip}$) $\to$ résiduel $\delta$ via sampler de Heun.}
  \caption{Architecture détaillée de \CorrDiff{}~\cite{mardani_corrdiff_2023}. Le premier étage régresse de manière déterministe la moyenne haute résolution $\muHR$ ; le second étage échantillonne par diffusion EDM le résiduel $\delta$ conditionnellement à l'étage 1.}
  \label{fig:corrdiff-archi}
\end{figure}

\subsection{Forces du paradigme à deux étages}
\label{ssec:soa-corrdiff-forces}

Trois forces structurelles caractérisent ce paradigme. \textbf{Stabilité d'entraînement} : la régression déterministe converge facilement, et la diffusion du second étage opère sur des résidus de distribution plus proche d'une gaussienne. \textbf{Calibration probabiliste explicite} : la variance est contrôlée par le second étage, indépendamment d'un équilibre adversarial. \textbf{Efficacité computationnelle} : l'échantillonnage par diffusion est limité au seul résiduel, de plus faible amplitude.

\subsection{Limites identifiées de CorrDiff}
\label{ssec:soa-corrdiff-limites}

L'examen critique fait apparaître deux limites principales. La \textbf{persistance de la sub-dispersion} : malgré la calibration explicite du second étage, CorrDiff reste sub-dispersif ($\sigma_{\mathrm{pred}}/\mathrm{RMSE}$ entre 0,3 et 0,5 selon les configurations), moins prononcée que les cGAN mais significative. L'\textbf{absence de structure causale explicite} : le premier étage est un U-Net générique apprenant $\muHR(\yLR, s_{\HR})$ sans contrainte sur la nature des dépendances ; rien ne distingue associations causales (humidité de basse couche $\to$ précipitation convective) et corrélations contingentes au régime d'entraînement. CorrDiff est donc vulnérable au changement de distribution, sans garantie d'intervention.

\subsection{Pourquoi CorrDiff plutôt que SR3 ?}
\label{ssec:soa-corrdiff-choix}

Le baseline non-causal d'\Oracle{} aurait pu être choisi parmi d'autres travaux de diffusion résiduelle. ResDiff et CorrDiff désignent la même filiation chez Mardani et~al. ; la question résiduelle est donc la comparaison à SR3~\cite{saharia_sr3_2022}. Trois raisons justifient CorrDiff. (i) \textbf{Décomposition formelle moyenne + résiduel} (équation~\ref{eq:corrdiff}) cohérente avec la statistique climatique : $\muHR$ capture le forçage-réponse, $\delta$ la variabilité fine stochastique ; SR3 ne propose pas cette séparation. (ii) \textbf{Conception pour le downscaling kilométrique climatique}, prise en compte de l'intermittence, queue lourde et hétérogénéité de la précipitation ; SR3 est une méthode générique. (iii) \textbf{Cadre architectural modulaire} : l'étage 1 peut être substitué sans toucher à l'étage 2 --- propriété qui permet précisément de \emph{causaliser} l'étage 1 (U-Net $\to$ module à graphe causal) en conservant l'étage 2 EDM intact.

% =============================================================
\section{Fonctions de perte et contraintes physiques}
\label{sec:soa-pertes}
% =============================================================

\subsection{Pertes adaptées à la précipitation}
\label{ssec:soa-pertes-precip}

Au-delà de la MSE, la littérature a développé plusieurs pertes adaptées à la précipitation. Les \textbf{pertes spectrales} (FACL, \emph{Fourier Amplitude Coherence Loss}) pénalisent l'écart entre spectres de Fourier prédit et cible, préservant le contenu fréquentiel fin et corrigeant partiellement le lissage MSE. La \textbf{régularisation Wasserstein} (héritée des WGAN) remplace la divergence Jensen-Shannon implicite par une distance Wasserstein-1, plus stable et mieux comportée aux extrêmes. Les \textbf{pondérations spécifiques aux extrêmes} (\emph{focal loss}, MSE pondérée par percentiles, top-$k$) accroissent le poids des pixels à forte intensité, concentrant l'apprentissage sur les cas extrêmes.

\subsection{Contraintes physiques et explicabilité}
\label{ssec:soa-physics-xai}

Des \textbf{contraintes physiques dures} (conservation de masse, monotonie $T_{\min} \leq T \leq T_{\max}$, positivité de la précipitation) sont imposées par construction architecturale ou projection en sortie ; elles améliorent la cohérence multivariée mais ne résolvent pas les verrous probabilistes. L'\textbf{explicabilité} (XAI) appliquée au downscaling (Grad-CAM, Integrated Gradients) révèle a posteriori les zones d'entrée influentes. Notre proposition emprunte une voie complémentaire : organiser l'architecture autour d'une structure causale apprise, apportant par construction cohérence et interprétabilité.

\begin{table}[ht]
  \centering
  \footnotesize
  \caption{Principales fonctions de perte et contraintes utilisées en downscaling climatique profond.}
  \label{tab:pertes}
  \begin{tabular}{p{3cm} p{3.5cm} p{5.5cm}}
    \toprule
    \textbf{Famille} & \textbf{Exemples} & \textbf{Cas d'usage principal} \\
    \midrule
    Pertes spectrales & FACL & Contenu fréquentiel fin\\
    Adversariales régularisées & Wasserstein, WGAN-GP & Stabilisation entraînement adversarial\\
    Pondération extrêmes & Focal, top-$k$, MSE pondérée & Pixels à forte intensité\\
    Contraintes dures & Conservation, monotonie, positivité & Cohérence physique multivariée\\
    Explicabilité a posteriori & Grad-CAM, Integrated Gradients & Diagnostic des dépendances\\
    Structure causale & DAG appris, intervention & Robustesse OOD et auditabilité (notre voie)\\
    \bottomrule
  \end{tabular}
\end{table}

Le tableau~\ref{tab:pertes} récapitule ces familles.

% =============================================================
\section{Cadres d'évaluation et foundation models}
\label{sec:soa-benchmarks}
% =============================================================

\subsection{Benchmarks standardisés}
\label{ssec:soa-benchmarks-std}

Trois benchmarks récents tentent de remédier à la fragmentation des protocoles d'évaluation. \textbf{WeatherBench~2}~\cite{rasp_weatherbench2_2023} étend le benchmark initial à la prévision météorologique pilotée par les données (RMSE, biais, ACC sur ERA5). \textbf{ClimateBench}~\cite{watson_climatebench_2022} cible la réponse à des scénarios d'émissions et l'émulation climatique. \textbf{ChaosBench}~\cite{nathaniel_chaosbench_2024} se concentre sur la prévision à grande échelle et la stabilité à long horizon. Le constat principal : \textbf{aucun benchmark public ne couvre le downscaling probabiliste à haute résolution} ; sa construction reste un travail ouvert.

\subsection{Foundation models climat}
\label{ssec:soa-foundation}

La période 2023-2024 a vu émerger des \emph{foundation models} climat. \textbf{ClimaX}~\cite{nguyen_climax_2023} (Microsoft) est un Transformer pré-entraîné sur CMIP6, capable de prévision, projection et downscaling après fine-tuning. \textbf{GraphCast}~\cite{lam_graphcast_2023} et \textbf{GenCast}~\cite{price_gencast_2024} (Google DeepMind) sont des graph neural networks entraînés sur ERA5 ayant dépassé plusieurs modèles physiques opérationnels. Le positionnement de notre travail est complémentaire : \Oracle{} cible le downscaling régional probabiliste, non le multi-tâche global ; la structure causale proposée pourrait à terme s'intégrer à une architecture de plus grande échelle.

% =============================================================
\section{Synthèse, tableau comparatif et gap scientifique}
\label{sec:soa-synthese}
% =============================================================

\subsection{Tableau comparatif modèle × verrou}
\label{ssec:soa-tableau}

Le tableau~\ref{tab:modele-verrou} récapitule les approches revues selon les trois axes du trilemme (stabilité, extrêmes, robustesse OOD), complétés par deux axes propres à notre contribution (structure causale explicite, caractère génératif probabiliste).

\begin{table}[ht]
  \centering
  \small
  \caption{Confrontation des principales approches du downscaling profond aux verrous du trilemme. Légende : \checkmark{} = bien traité, $\triangle$ = partiellement traité, $\times$ = non traité. La colonne \textbf{Causal} renvoie à l'indispensabilité structurelle (propriété~(O3)) qui sera formalisée au chapitre suivant.}
  \label{tab:modele-verrou}
  \begin{tabular}{l c c c c c c}
    \toprule
    \textbf{Approche} & \textbf{Année} & \textbf{Stabilité} & \textbf{Extrêmes} & \textbf{OOD} & \textbf{Causal} & \textbf{Génératif} \\
    \midrule
    DeepSD / U-Net & 2017 & \checkmark & $\times$ & $\times$ & $\times$ & $\times$ \\
    EDSR / SRCNN & 2017 & \checkmark & $\times$ & $\times$ & $\times$ & $\times$ \\
    FNO / SFNO & 2021-23 & \checkmark & $\triangle$ & $\triangle$ & $\times$ & $\times$ \\
    SwinIR / PrecipFormer & 2021-23 & \checkmark & $\triangle$ & $\times$ & $\times$ & $\times$ \\
    cGAN (Harris et al.) & 2022 & $\triangle$ & $\triangle$ & $\times$ & $\times$ & \checkmark \\
    Residual cGAN (Rampal) & 2024 & $\triangle$ & $\triangle$ & $\times$ & $\times$ & \checkmark \\
    VAE downscaling & 2022 & \checkmark & $\times$ & $\times$ & $\times$ & \checkmark \\
    Normalizing Flow & 2020 & \checkmark & $\triangle$ & $\times$ & $\times$ & \checkmark \\
    DDPM / EDM brut & 2020-22 & \checkmark & \checkmark & $\triangle$ & $\times$ & \checkmark \\
    SR3 / ResShift & 2022-23 & \checkmark & $\triangle$ & $\triangle$ & $\times$ & \checkmark \\
    \textbf{CorrDiff (Mardani)} & 2023 & \checkmark & $\triangle$ & $\triangle$ & $\times$ & \checkmark \\
    ClimaX (foundation) & 2023 & \checkmark & $\triangle$ & $\triangle$ & $\times$ & $\triangle$ \\
    \midrule
    \textbf{\Oracle{} (proposé)} & 2026 & \checkmark & $\triangle$ & $\triangle$ & \checkmark & \checkmark \\
    \bottomrule
  \end{tabular}
\end{table}

\noindent\emph{Note sur \Oracle{}.} La colonne « Causal » est attribuée à \Oracle{} au titre de la propriété (O3) --- architecturalement, la sortie ne peut pas contourner la matrice $\Adag$ apprise. Les triangles partiels en « Extrêmes » et « OOD » assument honnêtement les limites mesurées : RX1day mieux placé que la baseline mais F1@p99 régresse ; gain CRPS en inter-GCM historique mais pas de test sur scénario futur. Le détail est discuté au Chapitre~IV.

\subsection{Gap scientifique}
\label{ssec:soa-gap}

Le tableau~\ref{tab:modele-verrou} fait apparaître un constat structurant : \textbf{aucune des approches revues ne résout simultanément les trois axes du trilemme}. Les plus stables (U-Net, DDPM, EDM brut) ne garantissent pas la robustesse OOD ; les plus probabilistes (CorrDiff, cGAN résiduel) restent sub-dispersives sur les extrêmes ; les plus expressives (cGAN) souffrent d'instabilité. Aucune n'incorpore par construction de structure causale explicite : c'est précisément cette absence qui motive notre contribution.

Le gap scientifique se formule comme \textbf{la non-résolution du trilemme par les approches existantes}. Le paradigme à deux étages de \CorrDiff{} couvre le mieux le triplet stabilité + extrêmes + probabiliste, mais n'adresse pas la robustesse OOD par construction. Notre proposition vise à combler ce gap en \textbf{causalisant le premier étage} : remplacer le U-Net générique estimant $\muHR$ par un module structurellement causal contraignant cette estimation à passer par un graphe de dépendances appris.

\subsection{Positionnement d'Oracle dans le paysage}
\label{ssec:soa-oracle-position}

\begin{figure}[ht]
  \centering
  \figplaceholder{Figure F5 (Chap II) à produire --- cartouche conceptuel Oracle revisité avec références numériques : LR $\to$ Étage 1 causal (encodeur + DAG appris [Bello 2022, Pearl 2009]) $\to$ $\mu_{HR}$ ; bruit + $\mu_{HR}$ $\to$ Étage 2 diffusion EDM [Karras 2022] $\to$ $\delta$ ; sortie HR = baseline + $\mu_{HR}$ + $\delta$. Surligner « causalisation » comme contraste avec CorrDiff [Mardani 2023].}
  \caption{Cartouche conceptuel revisité du modèle \Oracle{} avec ancrage bibliographique. Le premier étage causalise \CorrDiff{}~\cite{mardani_corrdiff_2023} en y intégrant un graphe causal appris par contrainte DAGMA~\cite{bello_dagma_2022} ; le second étage conserve la diffusion EDM~\cite{karras_edm_2022}.}
  \label{fig:oracle-revisite}
\end{figure}

La figure~\ref{fig:oracle-revisite} résume la position de \Oracle{}. L'étage 1 est une causalisation de \CorrDiff{}~\cite{mardani_corrdiff_2023}, fondée sur la caractérisation différentiable de l'acyclicité de Bello et~al.~\cite{bello_dagma_2022} et le formalisme SCM de Pearl~\cite{pearl_causality_2009} ; l'étage 2 conserve la diffusion EDM~\cite{karras_edm_2022}. L'architecture détaillée fait l'objet du chapitre suivant.

\subsection{Transition vers le Chapitre III}
\label{ssec:soa-transition}

Le chapitre suivant détaille l'architecture du module causal récurrent (\emph{Recurrent Causal Network}, RCN) --- apprentissage simultané de la structure du graphe et de la dynamique de l'état latent, connexion au second étage de diffusion --- et explicite les objectifs O1-O6, en particulier la propriété d'intervention (O3) qui garantit par construction qu'une intervention sur la matrice d'adjacence apprise affecte la sortie.

% =============================================================
% Bibliographie
% =============================================================
\begin{thebibliography}{99}

\bibitem{dong_srcnn_2014} C.~Dong, C.~C.~Loy, K.~He, and X.~Tang. Learning a deep convolutional network for image super-resolution. In \emph{ECCV}, 2014.

\bibitem{lim_edsr_2017} B.~Lim, S.~Son, H.~Kim, S.~Nah, and K.~M.~Lee. Enhanced deep residual networks for single image super-resolution. In \emph{CVPR Workshops}, 2017.

\bibitem{vandal_deepsd_2017} T.~Vandal, E.~Kodra, S.~Ganguly, A.~Michaelis, R.~Nemani, and A.~R.~Ganguly. DeepSD: Generating high resolution climate change projections through single image super-resolution. In \emph{ACM SIGKDD}, 2017.

\bibitem{ronneberger_unet_2015} O.~Ronneberger, P.~Fischer, and T.~Brox. U-Net: Convolutional networks for biomedical image segmentation. In \emph{MICCAI}, 2015.

\bibitem{pan_pan_2019} J.~Pan et al. Spatial attention pyramid network for unsupervised single image super-resolution. \emph{arXiv preprint}, 2019.

\bibitem{li_fno_2021} Z.~Li, N.~Kovachki, K.~Azizzadenesheli, B.~Liu, K.~Bhattacharya, A.~Stuart, and A.~Anandkumar. Fourier neural operator for parametric partial differential equations. In \emph{ICLR}, 2021.

\bibitem{bonev_sfno_2023} B.~Bonev, T.~Kurth, C.~Hundt, J.~Pathak, M.~Baust, K.~Kashinath, and A.~Anandkumar. Spherical Fourier neural operators: Learning stable dynamics on the sphere. In \emph{ICML}, 2023.

\bibitem{liang_swinir_2021} J.~Liang, J.~Cao, G.~Sun, K.~Zhang, L.~Van~Gool, and R.~Timofte. SwinIR: Image restoration using Swin Transformer. In \emph{ICCV Workshops}, 2021.

\bibitem{lopezgomez_precipformer_2023} I.~Lopez-Gomez et al. PrecipFormer: A transformer-based architecture for precipitation downscaling. \emph{arXiv preprint}, 2023.

\bibitem{goodfellow_gan_2014} I.~Goodfellow, J.~Pouget-Abadie, M.~Mirza, B.~Xu, D.~Warde-Farley, S.~Ozair, A.~Courville, and Y.~Bengio. Generative adversarial nets. In \emph{NeurIPS}, 2014.

\bibitem{mirza_cgan_2014} M.~Mirza and S.~Osindero. Conditional generative adversarial nets. \emph{arXiv preprint arXiv:1411.1784}, 2014.

\bibitem{rampal_cgan_2024} N.~Rampal, P.~B.~Gibson, S.~Hobeichi, J.~Sherwood, and S.~B.~Power. Reliable generative downscaling of climate data with conditional GANs. \emph{Journal of Advances in Modeling Earth Systems}, 2024.

\bibitem{harris_2022} L.~Harris, A.~T.~T.~McRae, M.~Chantry, P.~D.~Dueben, and T.~N.~Palmer. A generative deep learning approach to stochastic downscaling of precipitation forecasts. \emph{Journal of Advances in Modeling Earth Systems}, 14(10):e2022MS003120, 2022.

\bibitem{price_rasp_2022} I.~Price and S.~Rasp. Increasing the accuracy and resolution of precipitation forecasts using deep generative models. In \emph{AISTATS}, 2022.

\bibitem{kingma_vae_2014} D.~P.~Kingma and M.~Welling. Auto-encoding variational Bayes. In \emph{ICLR}, 2014.

\bibitem{lange_vae_2022} S.~Lange et al. Variational autoencoders for climate downscaling. \emph{Environmental Data Science}, 2022.

\bibitem{groenke_flow_2020} B.~Groenke, L.~Madaus, and C.~Monteleoni. ClimAlign: Unsupervised statistical downscaling of climate variables via normalizing flows. \emph{arXiv preprint}, 2020.

\bibitem{ho_ddpm_2020} J.~Ho, A.~Jain, and P.~Abbeel. Denoising diffusion probabilistic models. In \emph{NeurIPS}, 2020.

\bibitem{karras_edm_2022} T.~Karras, M.~Aittala, T.~Aila, and S.~Laine. Elucidating the design space of diffusion-based generative models. In \emph{NeurIPS}, 2022.

\bibitem{saharia_sr3_2022} C.~Saharia, J.~Ho, W.~Chan, T.~Salimans, D.~J.~Fleet, and M.~Norouzi. Image super-resolution via iterative refinement. \emph{IEEE Transactions on Pattern Analysis and Machine Intelligence}, 2022.

\bibitem{yue_resshift_2023} Z.~Yue, J.~Wang, and C.~C.~Loy. ResShift: Efficient diffusion model for image super-resolution by residual shifting. In \emph{NeurIPS}, 2023.

\bibitem{lipman_flow_2023} Y.~Lipman, R.~T.~Q.~Chen, H.~Ben-Hamu, M.~Nickel, and M.~Le. Flow matching for generative modeling. In \emph{ICLR}, 2023.

\bibitem{song_consistency_2023} Y.~Song, P.~Dhariwal, M.~Chen, and I.~Sutskever. Consistency models. In \emph{ICML}, 2023.

\bibitem{rombach_latent_2022} R.~Rombach, A.~Blattmann, D.~Lorenz, P.~Esser, and B.~Ommer. High-resolution image synthesis with latent diffusion models. In \emph{CVPR}, 2022.

\bibitem{mardani_corrdiff_2023} M.~Mardani, N.~Brenowitz, Y.~Cohen, J.~Pathak, C.-Y.~Chen, C.-C.~Liu, A.~Vahdat, K.~Kashinath, J.~Kautz, and M.~Pritchard. Residual corrective diffusion modeling for km-scale atmospheric downscaling. \emph{arXiv preprint arXiv:2309.15214}, 2023.

\bibitem{rasp_weatherbench2_2023} S.~Rasp et al. WeatherBench 2: A benchmark for the next generation of data-driven global weather models. \emph{arXiv preprint arXiv:2308.15560}, 2023.

\bibitem{watson_climatebench_2022} D.~Watson-Parris et al. ClimateBench v1.0: A benchmark for data-driven climate projections. \emph{Journal of Advances in Modeling Earth Systems}, 2022.

\bibitem{nathaniel_chaosbench_2024} J.~Nathaniel et al. ChaosBench: A multi-channel, physics-based benchmark for subseasonal-to-seasonal climate prediction. \emph{NeurIPS Datasets and Benchmarks}, 2024.

\bibitem{nguyen_climax_2023} T.~Nguyen, J.~Brandstetter, A.~Kapoor, J.~K.~Gupta, and A.~Grover. ClimaX: A foundation model for weather and climate. In \emph{ICML}, 2023.

\bibitem{lam_graphcast_2023} R.~Lam et al. Learning skillful medium-range global weather forecasting. \emph{Science}, 382(6677):1416-1421, 2023.

\bibitem{price_gencast_2024} I.~Price et al. GenCast: Diffusion-based ensemble forecasting for medium-range weather. \emph{arXiv preprint}, 2024.

\bibitem{bello_dagma_2022} K.~Bello, B.~Aragam, and P.~Ravikumar. DAGMA: Learning DAGs via M-matrices and a log-determinant acyclicity characterization. In \emph{NeurIPS}, 2022.

\bibitem{pearl_causality_2009} J.~Pearl. \emph{Causality: Models, Reasoning, and Inference}. Cambridge University Press, 2nd edition, 2009.

\end{thebibliography}

\end{document}
